\documentclass{beamer}

\usepackage[utf8]{inputenc}
\usepackage[francais]{babel}
\usepackage{wrapfig}
\usepackage{amsmath}% http://ctan.org/pkg/amsmath
\usepackage{graphicx}

\usepackage{multirow}

\usepackage{sansmathaccent}
\pdfmapfile{+sansmathaccent.map}
\usepackage{tikz}
\usepackage{verbatim}

\setbeamertemplate{caption}{\raggedright\insertcaption\par}

\usetheme{Singapore}

\title{Détection d'anomalies de classification dans l'IoT via Machine Learning}
\author{Antoine Urban, Yohan Chalier}
\institute{Projet de filière SR2I \\ Télécom ParisTech}

\begin{document}

\begin{frame}
\titlepage
\end{frame}

\section{Introduction}

\begin{frame}{Introduction}
La détection d'obstacles: un enjeu de sécurité !
\begin{figure}
\centering
\includegraphics[width=.8\textwidth]{img/sensors.png}
\caption{Ensemble des capteurs présents dans le véhicule}
\end{figure}
\end{frame}

\begin{frame}{Attaques potentielles}

\begin{columns}

\begin{column}{0.5\textwidth}
\begin{figure}
\centering
\includegraphics[width=\textwidth]{img/blinding.png}
\caption{Attaque par aveuglement des capteurs}
\end{figure}
\end{column}

\begin{column}{0.5\textwidth}
\begin{figure}
\centering
\includegraphics[width=\textwidth]{img/data_flow.png}
\caption{Attaque par modification}
\end{figure}
\end{column}

\end{columns}

\end{frame}


\begin{frame}{Objectifs}
Proposition d'un modèle de classification multi-classes en réalisant un classeur à partir d'un algorithme d'apprentissage supervisé.
\end{frame}

\begin{frame}
\begin{figure}
\centering
\includegraphics[width=.9\textwidth]{img/structure.png}
\end{figure}
\end{frame}

\section{Démarche de recherche}

\begin{frame}{Première implémentation}
\begin{itemize}
\item Extraction des colonnes largeur et longueur de la base de données
\item Suppression des redondances
\item Définition de zones de décision arbitraires
\item Génération des données malicieuses
\end{itemize}
\begin{table}
\centering
\begin{tabular}{llll}
validité & intervalle de longueur & intervalle de largeur \\
\hline
non-malicieux & 3 à 6,5 mètres & 1,4 à 2,4 mètres \\
malicieux & 3 à 4,1 mètres & 2,05 à 2,4 mètres \\
malicieux & 5,25 à 6,5 mètres & 1,4 à 1,65 mètres \\
\end{tabular}
\end{table}
\end{frame}

\begin{frame}
\begin{figure}
\centering
\includegraphics[width=.8\textwidth]{img/first_try.png}
\end{figure}
\end{frame}

\begin{frame}{Méthodes d'évaluation}
\framesubtitle{Matrice de confusion}

\begin{tabular}{ l | l | c | c | c c}
\multicolumn{2}{c}{} & \multicolumn{2}{c}{Classe réelle} & \\
\cline{3-4}
\multicolumn{2}{c|}{} & Positif & Négatif \\
\cline{2-4}
\multirow{2}{*}{Classe prédite} & Positif & $TP$ & $FP$ & $PPV$ & $FDR$ \\
\cline{2-4}
& Négatif & $FN$ & $TN$ & $FOR$ & $NPV$ \\
\cline{2-4}
\multicolumn{1}{r}{} & \multicolumn{1}{l}{} & \multicolumn{1}{c}{$TPR$} & \multicolumn{1}{c}{$FPR$} \\
\multicolumn{1}{l}{} & \multicolumn{1}{l}{} & \multicolumn{1}{c}{$FNR$} & \multicolumn{1}{c}{$TNR$} \\
\end{tabular}

\end{frame}

\begin{frame}{Méthodes d'évaluation}
\framesubtitle{Score F1}

\begin{block}{Objectif}
Maximisation du score F1 comme critère de performance
\end{block}

\begin{equation}
\text{f1-score} = \dfrac{2\times(\text{Recall} \times \text{Precision})}{(\text{Recall} + \text{Precision})} = 2\times\dfrac{PPV \times TPR}{PPV + TPR}
\end{equation}

\begin{equation}
\text{Precision} = \dfrac{TP}{TP+FP}
\end{equation}

\begin{equation}
\text{Recall} = \dfrac{TP}{TP+FN}
\end{equation}


\end{frame}

\end{document}
